% ============================================
% Заключение
% ============================================

\chapter*{ЗАКЛЮЧЕНИЕ}
\addcontentsline{toc}{chapter}{ЗАКЛЮЧЕНИЕ}

\section*{Итоги работы}
\addcontentsline{toc}{section}{Итоги работы}

В ходе выполнения настоящей научно-исследовательской работы был проведён комплексный анализ предметной области на стыке системной динамики и больших языковых моделей, спроектирована и реализована программная платформа \texttt{sdCloud}, предназначенная для автоматизированной генерации системно-динамических моделей на основе текстовых описаний.

В рамках работы были достигнуты следующие результаты:

\begin{enumerate}
    \item \textbf{Обзор предметной области.} Проведён систематический анализ теоретических основ системной динамики, включая концептуальное моделирование (CLD), построение диаграмм потоков и уровней (SFD) и формализацию уравнений. Исследованы существующие программные платформы моделирования (Vensim, Stella, AnyLogic, PySD), технологии больших языковых моделей (GPT-4, Claude, семейства LLaMA и Mistral) и существующие работы по интеграции LLM с системной динамикой (ChatPySD, исследования Akhavan, Liu и др.).

    \item \textbf{Проектирование архитектуры.} Разработана архитектура платформы \texttt{sdCloud}, включающая backend на базе FastAPI с использованием SQLAlchemy ORM, PostgreSQL, Redis и асинхронной очереди задач Taskiq; frontend на базе Next.js с применением методологии Feature-Sliced Design; а также модули генерации моделей, валидации графов и интеграции с OpenRouter API.

    \item \textbf{Разработка конвейера генерации.} Предложен и реализован многоэтапный конвейер генерации системно-динамических моделей, включающий извлечение переменных, построение CLD, генерацию уравнений, построение SFD, детерминированную генерацию Python-кода и Vensim-экспорт. Конвейер использует few-shot промптинг для структурированной генерации данных LLM.

    \item \textbf{Детерминированный генератор Python-кода.} Разработан модуль \texttt{python\_generator.py}, формирующий исполняемый код по структурированным уравнениям без повторного обращения к LLM, что повышает воспроизводимость и снижает влияние галлюцинаций.

    \item \textbf{Fallback-симуляция через SciPy.} Реализован модуль \texttt{scipy\_runner.py}, обеспечивающий симуляцию сгенерированных моделей методами \texttt{solve\_ivp} при невозможности прямой загрузки в PySD.

    \item \textbf{Реализация модуля валидации.} Разработан модуль валидации графовых структур на базе библиотеки NetworkX, выполняющий проверку связности графа, обнаружение петель обратной связи, верификацию структурных элементов SFD и согласованность переменных, уравнений и CLD.

    \item \textbf{Гостевой режим.} Реализована возможность анонимной генерации моделей и симуляций с сохранением истории в браузере, что снижает барьер входа для новых пользователей.

    \item \textbf{Программная реализация.} Backend-приложение реализовано на Python с использованием FastAPI, frontend-приложение --- на TypeScript с использованием Next.js и Bun. Система включает REST API с эндпоинтами для авторизации, диалогов, моделей, симуляций и гостевого режима, реляционную схему базы данных, модульные и интеграционные тесты.

    \item \textbf{Тестирование.} Проведено тестирование платформы на типовых задачах системной динамики, подтвердившее работоспособность конвейера генерации и симуляции для моделей низкой и средней сложности. Выявлены ограничения бесплатных 30B-моделей на OpenRouter.
\end{enumerate}

\section*{Научная новизна}
\addcontentsline{toc}{section}{Научная новизна}

Научная новизна полученных результатов заключается в следующем:

\begin{enumerate}
    \item Предложен многоэтапный конвейер генерации системно-динамических моделей с промежуточными представлениями в виде CLD и SFD, обеспечивающий прозрачность и интерпретируемость процесса генерации.

    \item Разработан метод валидации графовых структур системно-динамических моделей, основанный на анализе топологических свойств графа причинно-следственных связей. Метод включает проверку связности, обнаружение петель обратной связи (алгоритм Джонсона) и верификацию структурных ограничений SFD.

    \item Предложен детерминированный метод генерации исполняемого Python-кода по структурированным уравнениям, заменяющий нестабильную генерацию кода LLM и повышающий воспроизводимость симуляций.

    \item Разработан fallback-механизм симуляции на основе SciPy \texttt{solve\_ivp}, обеспечивающий работоспособность платформы при несовместимости сгенерированного кода с текущей версией PySD.

    \item Создана программная платформа \texttt{sdCloud}, представляющая собой комплексное решение, интегрирующее веб-технологии, базы данных, асинхронную обработку задач, API больших языковых моделей и гостевой режим в единую систему для AI-ассистированного системно-динамического моделирования.
\end{enumerate}

\section*{Практическая значимость}
\addcontentsline{toc}{section}{Практическая значимость}

Практическая значимость работы определяется следующими факторами:

\begin{enumerate}
    \item \textbf{Снижение барьера входа.} Платформа \texttt{sdCloud} позволяет специалистам без глубокой подготовки в системной динамике создавать и анализировать системно-динамические модели, что расширяет круг потенциальных пользователей методологии.

    \item \textbf{Ускорение процесса моделирования.} Автоматическая генерация моделей из текстовых описаний сокращает время создания прототипа модели с часов до минут.

    \item \textbf{Образовательный потенциал.} Платформа может использоваться в учебном процессе для демонстрации принципов системной динамики и проверки понимания студентами.

    \item \textbf{Открытый исходный код.} Программная реализация платформы доступна в виде открытого исходного кода, что позволяет сообществу развивать и адаптировать систему под специфические задачи.

    \item \textbf{Архитектурные решения.} Применённые в проекте архитектурные паттерны (многоуровневая архитектура, FSD, асинхронная обработка задач) могут быть использованы при разработке аналогичных систем в смежных предметных областях.
\end{enumerate}

\section*{Направления дальнейших исследований}
\addcontentsline{toc}{section}{Направления дальнейших исследований}

На основе проведённого исследования выделены следующие перспективные направления дальнейшей работы:

\begin{enumerate}
    \item Улучшение визуализации графов (автоматическая раскладка, масштабирование текста) и добавление редактора моделей в интерфейс.

    \item Интеграция более мощных и стабильных языковых моделей (Claude 3.5 Sonnet, GPT-4o-mini, Gemini Flash) для повышения качества сложной генерации.

    \item Исследование возможностей донастройки (fine-tuning) открытых LLM на специализированном корпусе данных системной динамики.

    \item Разработка методов иерархической декомпозиции сложных систем и анализа чувствительности.

    \item Внедрение механизмов обратной связи от пользователя для итеративного уточнения и коррекции сгенерированных моделей.

    \item Переход к мультиагентной архитектуре генерации после обеспечения надёжности structured output и стабильности используемых LLM.

    \item Расширение платформы функциями коллаборативного моделирования, публикации моделей и интеграции с историческими данными для калибровки параметров.

    \item Проведение комплексного пользовательского исследования (user study) для оценки практической полезности платформы.
\end{enumerate}

Таким образом, выполненная работа создаёт основу для дальнейших исследований и разработок на стыке искусственного интеллекта и системной динамики, открывая новые возможности для democratизации методологии моделирования сложных систем.
